%%%%%%%%%%  scan idx 403  |  folio 390  %%%%%%%%%%
\SGAPrintPage{390}%
% SGA_NATIVE_CONTINUATION_389_390
\providecommand{\Sym}{\operatorname{Sym}}
\providecommand{\GL}{\operatorname{GL}}
\providecommand{\Grass}{\operatorname{Grass}}
\providecommand{\gr}{\operatorname{gr}}
\providecommand{\ctop}{c_{\mathrm{top}}}
\providecommand{\Rep}{\operatorname{Rep}}
\providecommand{\RZ}{R_{\mathbb Z}}
\providecommand{\Dflag}{D}
\providecommand{\Gmn}{G}
complete the induction, noting that, if $0\leq s\leq r$, and if $x_j=0$ for $s<j\leq r$, we have:
\[
  \sum_{i=0}^{s}f^*(x_i)\xi^i=0,
\]
whence:
\[
  \sum_{i=0}^{s}f^*(x_i)\xi^{i+r-s}=0,
\]
whence $x_s=0$ by what precedes.

\sgaStableTarget{sga6.e06.sec.3}\section*{3. Consequences of the structure theorem for the \(K^{\bullet}\) of a projective bundle}

\sgaStableTarget{sga6.e06.prop.3-1}\textbf{Proposition 3.1.} \uline{Let} \(S\) \uline{be a quasi}-\uline{compact scheme}, \(\underline{E}_1,\ldots,\underline{E}_n\) \uline{locally free} \(\underline{O}_S\)\uline{-Modules of finite type}, \(\mathbf D(\underline{E}_i)\) \uline{the flag scheme of} \(\underline{E}_i\), \(X=\mathbf D(\underline{E}_1)\times_S\ldots\times_S \mathbf D(\underline{E}_n)\).
\uline{Then the canonical homomorphism}:
\[
  K^{\bullet}(S)\longrightarrow K^{\bullet}(X)
\]
\uline{is injective} (\(K^{\bullet}\) \uline{being defined}, \uline{either from locally free} \(\underline{O}_S\)\uline{-Modules of finite type}, \uline{or from perfect complexes}).

We check immediately, using the universal property of flag schemes, that \(X\) can be constructed as follows: we consider the sequence of \(S\)-schemes \(X_1,\ldots,X_n\), defined inductively by taking \(X_k\) to be the flag scheme of the inverse image of \(\underline{E}_k\) over \(X_{k-1}\); then \(X=X_n\). We then see at once that, to prove the proposition, it suffices to prove it in the case \(n=1\), the general case following by induction. We can therefore assume that \(X\) is the flag scheme of a locally free \(\underline{O}_S\)-Module of finite type \(\underline{E}\).

We know that \(X\) can then be constructed as follows: we put \(\underline{E}=\underline{F}_0\), \(S=X_0\), and define inductively schemes \(X_1,\ldots,X_n\) and
%%%%%%%%%%  scan idx 404  |  folio 391  %%%%%%%%%%
\SGAPrintPage{391}%
for every \(i\) a locally free \(\underline{O}_{X_i}\)-Module of finite type \(\underline{F}_i\) by putting:
\[
  X_k=\mathbf P(\underline{F}_{k-1});\qquad
  \underline{F}_k=\Ker\bigl(f_k^*(\underline{F}_{k-1})\longrightarrow \underline{O}_{X_k}(1)\bigr),
\]
where \(f_k\) is the structural morphism \(X_k\to X_{k-1}\).

We then see that it suffices to prove the proposition in the case where \(X\) is a projective bundle over \(S\), the case of a flag bundle following by induction. This case follows immediately from 1.1 (resp. 2.1).

\sgaStableTarget{sga6.e06.thm.3-2}\textbf{Theorem 3.2.} \uline{Let} \(S\) \uline{be a quasi-compact scheme}. \uline{The pre}-\(\lambda\)-\uline{ring} \(K^{\bullet}(S)\), \uline{defined from locally free} \(\underline{O}_S\)\uline{-Modules of finite type}, \uline{is a} \(\lambda\)-\uline{ring}~(V 2.4).

The homomorphism \(\lambda_t:K^{\bullet}(S)\longrightarrow 1+\widehat{K^{\bullet}(S)[[t]]}^{+}\) being additive, it suffices, to check that it is a \(\lambda\)-homomorphism, to make the verification for elements of the form \(\cl(\underline{E})\), where \(\underline{E}\) is a locally free \(\underline{O}_S\)-Module of finite type. We are therefore reduced to showing that, for every pair of elements \(E,F\) of \(K^{\bullet}(S)\), classes of locally free \(\underline{O}_S\)-Modules of finite type \(\underline{E},\underline{F}\), relations V 2.4.2 hold:
\[
  \lambda^i(EF)=P_i\bigl(\lambda^1(E),\ldots,\lambda^i(E);\lambda^1(F),\ldots,\lambda^i(F)\bigr);
\]
\[
  \lambda^j(\lambda^i(E))=Q_{i,j}\bigl(\lambda^1(E),\ldots,\lambda^{ij}(E)\bigr).
\]

Since \(S\) is a quasi-compact scheme, we can find a finite covering of \(S\) by disjoint opens such that, on each, \(\underline{E}\) and \(\underline{F}\) have constant rank. The \(K^{\bullet}\) of a finite direct sum of opens is the product of the \(K^{\bullet}\) of these opens. It therefore suffices to check relations V 2.4.2 in each of the \(K^{\bullet}\) of the opens of the covering. We can therefore assume \(\underline{E}\) and \(\underline{F}\) to have constant rank on \(S\).

%%%%%%%%%%  scan idx 405  |  folio 392  %%%%%%%%%%
\SGAPrintPage{392}%
Now let \(X=\mathbf D(\underline{E})\times_S \mathbf D(\underline{F})\). The \(\lambda\)-homomorphism \(K^{\bullet}(S)\to K^{\bullet}(X)\) is injective by 3.1. Consequently, to prove relations V 2.4.2, we can work in \(K^{\bullet}(X)\). Since in \(K^{\bullet}(X)\) the elements \(E\) and \(F\) decompose into sums of classes of invertible sheaves, it suffices, by additivity, to prove these same relations for elements \(\xi,\eta\), classes of invertible sheaves; they can also be written:
\[
  1+\xi\eta t=(1+\xi t)\circ(1+\eta t);
  \qquad
  \lambda_u(1+\xi t)=1+\{1+\xi t\}u,
\]
by putting them in the form V 2.4.1. They are then precisely relations V 2.3.1 and 2.3.3 defining the \(\lambda\)-ring structure of \(1+\widehat{K^{\bullet}(S)[[t]]}^{+}\).

Let us state the following more general result:

\sgaStableTarget{sga6.e06.thm.3-3}\textbf{Theorem 3.3.} \uline{Let} \((T,\underline{A})\) \uline{be a commutatively ringed topos}; \uline{write} \(K\) \uline{for the Grothendieck group of} \(\underline{A}\)-\uline{Modules locally free} of bounded rank on \(T\). \uline{Then}, \(K\), \uline{endowed with its canonical pre}-\(\lambda\)-\uline{structure}~(V 2.2 b)), \uline{is} \uline{a} \(\lambda\)-\uline{ring}.

To show that \(K\) is a \(\lambda\)-ring, we must prove relations V 2.4.1, and, by additivity, it suffices to prove them for elements \(E,F\) which are classes of locally free \(\underline{A}\)-Modules of bounded rank. Since the ranks of \(\underline{E}\) and \(\underline{F}\) are bounded, we can decompose \(T\) into a finite ``disjoint sum'' of opens over which \(\underline{E}\) and \(\underline{F}\) have constant rank. Using the commutation of \(K^{\bullet}\) with finite direct sums, and putting relations V 2.4.1 in the form V 2.4.2, we see that it suffices to prove relations V 2.4.2 in each of the \(K^{\bullet}\) of the opens considered, that is, it suffices to prove relations V 2.4.2 for elements \(E,F\), classes of \(\underline{A}\)-Modules \(\underline{E},\underline{F}\) locally free of constant rank.

%%%%%%%%%%  scan idx 406  |  folio 393  %%%%%%%%%%
\SGAPrintPage{393}%
Let \(\underline{E}\) be a locally free \(\underline{A}\)-Module of finite rank \(n\). There exists a torsor \(\underline{P}\) under \(\underline{\GL}(n,\underline{A})\) such that \(\underline{E}=\underline{P}\times_{\underline{\GL}(n,\underline{A})} (\underline{A})^n\), where \((\underline{A})^n\) is regarded as a \(\underline{\GL}(n,\underline{A})\)-Module by the standard representation. We can then define a \(\lambda\)-homomorphism from \(R_{\mathbb Z}(\mathbf{Gl}(n))\) to \(K\) by associating to every representation \(\rho\) of \(\mathbf{Gl}(n)\) in a free \(\mathbb Z\)-module of rank \(p\) the class of the locally free \(\underline{A}\)-Module
\(\underline{P}\times_{\underline{\GL}(n,\underline{A})} (\underline{A})^p\), where \((\underline{A})^p\) is regarded as a \(\underline{\GL}(n,\underline{A})\)-module by the representation \(\rho\). Under this homomorphism, the class of \(\underline{E}\) is therefore the image of the standard representation. To check a polynomial relation among the \(\lambda^i(E)\), it therefore suffices that this relation hold for the standard representation in \(R_{\mathbb Z}(\mathbf{Gl}(n))\).

We can argue similarly with two locally free \(\underline{A}\)-Modules \(\underline{E},\underline{F}\) of ranks \(n,p\); we can then find a torsor \(\underline{P}\) under \(\underline{\GL}(n,\underline{A})\times\underline{\GL}(p,\underline{A})\) such that \(\underline{E}=\underline{P}\times_{\underline{\GL}(n,\underline{A})\times\underline{\GL}(p,\underline{A})} (\underline{A})^n\), where \((\underline{A})^n\) is regarded as a \(\underline{\GL}(n,\underline{A})\times\underline{\GL}(p,\underline{A})\)-Module by letting \(\underline{\GL}(n,\underline{A})\) act by the standard representation and \(\underline{\GL}(p,\underline{A})\) by the trivial representation, and conversely for \(\underline{F}\). We then define a \(\lambda\)-homomorphism from \(R_{\mathbb Z}(\mathbf{Gl}(n)\times\mathbf{Gl}(p))\) to \(K\) in the same way as above, and are thus reduced to checking the relations among the \(\lambda^i(E)\), \(\lambda^j(F)\) for the corresponding elements of \(R_{\mathbb Z}(\mathbf{Gl}(n)\times\mathbf{Gl}(p))\).

We are therefore reduced to showing that, for every group scheme \(G\) over \(\Spec\mathbb Z\) which is a finite product of linear group schemes, the pre-\(\lambda\)-ring \(R_{\mathbb Z}(G)\) is a \(\lambda\)-ring. For this, we can refer to [5], where \(R_{\mathbb Z}(G)\) is computed explicitly. To avoid using this result, we can reason in a manner analogous to 3.2, using 1.14.

